\documentclass[11pt,a4paper]{article}

\usepackage[utf8]{inputenc}
\usepackage[T1]{fontenc}
\usepackage[spanish,es-noquoting]{babel}
\usepackage[margin=2.5cm]{geometry}
\usepackage{amsmath}
\usepackage{booktabs}
\usepackage{hyperref}

\hypersetup{colorlinks=true, linkcolor=black, urlcolor=blue}

\title{Modelo termohidráulico de PPHE por la ecuación (3)\\
       \large Contraste contra el segundo caso de estudio de O. Arsenyeva et al.}
\author{Informe de cambios en el código}
\date{15 de septiembre de 2026}

\begin{document}
\maketitle

\section{Objetivo}

Se reorienta el modelo para validarlo contra el \textbf{segundo} caso de estudio del artículo
de Arsenyeva et al.\footnote{O. Arsenyeva, J. Tran, M. Piper, E. Kenig, \emph{An approach for
pillow plate heat exchangers design for single-phase applications}, Appl. Therm. Eng. 147
(2019) 579--591.} (tren de precalentamiento de crudo), en lugar del primero. El encargo
concreto fue: tomar el número de chapas y la separación entre paneles de la Tabla 9 como datos
de entrada, y obtener las velocidades con la ecuación (3) del artículo,
\begin{equation}
  w = \frac{G}{\rho \cdot f_{ch}},
\end{equation}
eliminando el método iterativo sobre la ecuación implícita (28).

Además se restauran en la geometría los valores de $b$ y $b_i$ del artículo.

\section{Resumen de cambios}

\begin{table}[h]
\centering
\begin{tabular}{llp{8.2cm}}
\toprule
Fichero & Estado & Cambio \\
\midrule
\texttt{geometría.py}      & Modificado & $b$ restaurado a los valores del artículo; comentario aclarando su procedencia \\
\texttt{formulas.py}       & Modificado & Cuatro funciones nuevas: ecuaciones (16), (17), (3) y (19) \\
\texttt{modelo\_antiguo.py} & Reescrito  & Modelo completo por la vía de la ecuación (3), contrastado contra la Tabla 9 \\
\bottomrule
\end{tabular}
\end{table}

\subsection{\texttt{geometría.py}}

Los valores de $b$ estaban en 3, 2 y 2\,mm por pruebas anteriores. Se restauran a 5.5, 7.5 y
20\,mm. Los de $b_i$ (3.4, 3.0 y 7.0\,mm) ya eran correctos.

Se añade un comentario al campo porque su procedencia es fácil de confundir: \textbf{$b$ no
aparece en la Tabla 1} del artículo, que solo recoge la geometría del panel
($\delta_{pp}$, $b_i$, $2s_L$, $s_T$, $w_{pp}$, $L_{pp}$, $w_e$). En el método del artículo $b$
es la \emph{variable de diseño} que se barre; los tres valores anteriores son los óptimos que
reporta su Tabla 5 para el primer caso de estudio.

\subsection{\texttt{formulas.py}}

Se añaden las ecuaciones que faltaban para poder trabajar desde el caudal:

\begin{align}
  f_{chI} &= \frac{b_i}{\sqrt{2}}\,(w_{pp}-2w_e) & &\text{(16)}\\
  f_{chE} &= \left[(b_i+b)-\frac{b_i}{\sqrt{2}}\right](w_{pp}-2w_e) & &\text{(17)}\\
  w &= \frac{G}{\rho\,N\,f_{ch}} & &\text{(3)}\\
  \Delta P &= \zeta\,\frac{L_F}{d_e}\,\frac{\rho w^2}{2} + \zeta_{DZ}\,\rho w^2 & &\text{(19)}
\end{align}

La ecuación (19) es la inversa de la ya existente \texttt{LF\_ploss}: aquella despeja la
longitud conocida la pérdida de carga, y ésta la pérdida de carga conocida la longitud.
Nótese que $f_{ch} = d_e/2 \cdot (w_{pp}-2w_e)$ en ambos canales, es decir, el hueco medio por
la anchura útil.

\subsection{\texttt{modelo\_antiguo.py}}

Se reescribe por completo. Desaparece la búsqueda iterativa del número de canales y la
estructura pasa a recorrer las seis variantes de la Tabla 9 (todas ellas geometría PPHE2, con
$b$ distinto en cada una). Se conserva la estructura de \texttt{modelo\_termohidraulico.py}:
\emph{dataclass} de resultados, función de impresión, diccionarios de propiedades y bucle con
los bloques \textsc{velocidades} / \textsc{cálculo interno} / \textsc{cálculo externo}.

Se añade una segunda tabla de salida que contrasta lo calculado con lo que reporta la Tabla 9.

\section{Dos interpretaciones que ha habido que deducir}

El artículo no define con claridad dos cosas necesarias para el cálculo. Ambas se han resuelto
por consistencia numérica con sus propios datos.

\subsection{La columna $n_{pl}$ son chapas, no canales}

Si se toma $n_{pl}$ como número de canales internos, las densidades que exige la Tabla 9 al
aplicar la ecuación (3) resultan imposibles para hidrocarburos líquidos. Tomando
$N = n_{pl}/2$ salen valores normales y, sobre todo, \emph{decrecen de forma monótona con la
temperatura media de la corriente}, que es la comprobación que de verdad discrimina.

\begin{table}[h]
\centering
\begin{tabular}{crrrrrr}
\toprule
& \multicolumn{2}{c}{$N = n_{pl}$} & \multicolumn{2}{c}{$N = n_{pl}/2$} & \multicolumn{2}{c}{$T$ media [\textdegree C]}\\
\cmidrule(lr){2-3}\cmidrule(lr){4-5}\cmidrule(lr){6-7}
HE & $\rho_{int}$ & $\rho_{ext}$ & $\rho_{int}$ & $\rho_{ext}$ & int & ext\\
\midrule
1 & 384.8 & 360.2 & 769.6 & 720.4 & 166 & 154\\
2 & 355.4 & 390.6 & 710.8 & 781.1 & 187 & 100\\
3 & 311.3 & 327.7 & 622.5 & 655.4 & 243 & 170\\
4 & 256.0 & 287.9 & 511.9 & 575.8 & 280 & 238\\
5 & 274.4 & 308.7 & 548.8 & 617.4 & 277 & 236\\
6 & 253.0 & 274.5 & 505.9 & 548.9 & 323 & 277\\
\bottomrule
\end{tabular}
\caption{Densidades en kg/m$^3$ que implica la Tabla 9 bajo cada lectura de $n_{pl}$.}
\end{table}

La lectura encaja con la Nomenclatura del artículo, donde $N$ son los \emph{paneles} y el
subíndice $pl$ denota \emph{chapa}, y explica que la ecuación (13) obligue a que $n_{pl}$ sea
par: cada panel se forma soldando dos chapas.

\subsection{El área de intercambio se calcula con las chapas}

Como comprobación independiente de lo anterior, el área de la Tabla 9 se reproduce con
$F = n_{pl}\,L_F\,(w_{pp}-2w_e)\,F_x$, usando el número de chapas. Con el número de canales
aparece un factor $2$ sistemático.

\begin{table}[h]
\centering
\begin{tabular}{crrrrr}
\toprule
HE & $n_{pl}L_F W F_x$ & $F$ (Tabla 9) & cociente & $N L_F W F_x$ & cociente\\
\midrule
1 & 14.68 & 15.85 & 1.080 & 7.34 & 2.159\\
2 &  8.81 &  9.43 & 1.070 & 4.40 & 2.141\\
3 &  9.08 &  9.09 & 1.001 & 4.54 & 2.002\\
4 & 13.09 & 12.71 & 0.971 & 6.54 & 1.942\\
5 & 18.54 & 19.13 & 1.032 & 9.27 & 2.063\\
6 & 12.07 & 11.79 & 0.977 & 6.04 & 1.953\\
\bottomrule
\end{tabular}
\caption{Área de intercambio en m$^2$, calculada con la propia $L_F$ del artículo.}
\end{table}

Es coherente con el factor 2 de la ecuación (24): cada canal interno está envuelto por dos
chapas, de modo que el área total es $2\,N\,F_{pl} = n_{pl}\,F_{pl}$. Ambas lecturas se
sostienen mutuamente.

\section{Verificación previa: la Tabla 9 es coherente}

En trabajo anterior se determinó que las velocidades de la Tabla 5 (primer caso de estudio) no
son reproducibles desde los datos de partida del propio artículo, por estar calculadas con el
cociente de secciones de la ecuación (30) invertido. Antes de usar la Tabla 9 como referencia
se comprobó que no arrastra el mismo problema.

Partiendo de $w = G/(\rho N f_{ch})$ en los dos canales, el cociente de velocidades fija la
relación de densidades. Solo una de las dos versiones da valores posibles:

\begin{table}[h]
\centering
\begin{tabular}{crrrrr}
\toprule
HE & $b$ [mm] & $f_{chE}/f_{chI}$ & $w_{ext}/w_{int}$ & $\rho_{int}/\rho_{ext}$ correcta & invertida\\
\midrule
1 & 1.5 & 1.1213 & 2.7638 & 1.068 & 0.850\\
2 & 1.0 & 0.8856 & 2.8134 & 0.910 & 1.160\\
3 & 2.0 & 1.3570 & 2.8537 & 0.950 & 0.516\\
4 & 0.0 & 0.4142 & 1.7271 & 0.889 & 5.182\\
5 & 0.5 & 0.6499 & 1.6123 & 0.889 & 2.104\\
6 & 1.0 & 0.8856 & 3.5193 & 0.922 & 1.175\\
\bottomrule
\end{tabular}
\end{table}

La versión correcta da relaciones entre 0.889 y 1.068, todas plausibles. La invertida exigiría
que en HE4 el producto caliente fuese 5.18 veces más denso que el crudo. El argumento decisivo
son las variantes 4 y 5, que comparten fluidos (Producto 3 y 4 a 289.9\,\textdegree C contra
crudo a 224.6\,\textdegree C): la versión correcta da \emph{exactamente} el mismo 0.889 en
ambas, mientras que la invertida da 5.182 y 2.104.

El motivo de la diferencia con el primer caso de estudio es claro: aquí el fluido caliente
tiene menor caudal y va al canal interno, que es la configuración para la que se dedujeron las
ecuaciones (22)--(28). En ella los dos significados del subíndice 1 --- corriente caliente y
canal I --- coinciden, y no hay ambigüedad posible. \textbf{La Tabla 9 es, por tanto, mejor
banco de pruebas que la Tabla 5.}

\section{Resultados}

\subsection{Caso de estudio 2 frente a la Tabla 9}

\begin{table}[h]
\centering
\small
\begin{tabular}{crrrrrrrrrrrr}
\toprule
& \multicolumn{2}{c}{$u_{int}$ [m/s]} & \multicolumn{2}{c}{$u_{ext}$ [m/s]} & \multicolumn{2}{c}{$L_F$ [m]} & \multicolumn{2}{c}{$U$ [W/m$^2$K]} & \multicolumn{2}{c}{$F$ [m$^2$]} & \multicolumn{2}{c}{$\Delta P_{int}$ [kPa]}\\
\cmidrule(lr){2-3}\cmidrule(lr){4-5}\cmidrule(lr){6-7}\cmidrule(lr){8-9}\cmidrule(lr){10-11}\cmidrule(lr){12-13}
HE & calc & T9 & calc & T9 & calc & T9 & calc & T9 & calc & T9 & calc & T9\\
\midrule
1 & 1.943 & 1.943 & 5.370 & 5.370 & 2.411 & 2.250 & 2986 & 2863 & 15.73 & 15.85 & 155.5 & 100\\
2 & 2.229 & 2.229 & 6.271 & 6.271 & 1.194 & 1.350 & 3261 & 2969 &  7.79 &  9.43 &  95.2 & 100\\
3 & 2.051 & 2.051 & 5.853 & 5.853 & 1.063 & 1.670 & 2754 & 2125 &  5.78 &  9.09 &  65.9 & 100\\
4 & 3.346 & 3.346 & 5.779 & 5.779 & 0.615 & 0.830 & 3269 & 3316 &  9.70 & 12.71 &  82.6 & 100\\
5 & 2.920 & 2.920 & 4.708 & 4.708 & 0.841 & 1.100 & 2990 & 2937 & 14.17 & 19.13 &  90.7 & 100\\
6 & 1.947 & 1.947 & 6.853 & 6.852 & 1.025 & 1.850 & 2510 & 3050 &  6.69 & 11.79 &  49.0 & 100\\
\bottomrule
\end{tabular}
\end{table}

Lectura de la tabla, con cuidado de no atribuirle más valor del que tiene:

\begin{itemize}
  \item Las \textbf{velocidades coinciden exactamente}, pero la comparación es
        \emph{circular}: las densidades se han despejado de la propia Tabla 9 a falta de las
        reales. Lo que sí confirma es que la lectura $N = n_{pl}/2$ es autoconsistente.
  \item El \textbf{área de HE1} sale 15.73 frente a 15.85\,m$^2$, un 0.8\,\% de error, con la
        corrección del apartado anterior.
  \item $U$ cae dentro del 10--30\,\% pese a usar propiedades genéricas, lo que indica que el
        orden de magnitud es razonable, pero no constituye validación.
  \item $L_F$, $F$ y las pérdidas de carga arrastran el error de las propiedades provisionales.
\end{itemize}

\subsection{Caso de estudio 1 (\texttt{modelo\_termohidraulico.py})}

Sin cambios en este informe; se recoge el estado actual como referencia.

\begin{table}[h]
\centering
\begin{tabular}{crrrrrrrr}
\toprule
Caso & $u_{int}$ & $u_{ext}$ & $Re_{int}$ & $Re_{ext}$ & $h_{int}$ & $h_{ext}$ & $U$ & $L_F$\\
\midrule
PPHE1 & 1.3591 & 0.6707 & 8.01e3 & 1.71e4 & 6556.96 & 6414.03 & 2790.04 & 3.8474\\
PPHE2 & 0.8396 & 0.2834 & 4.36e3 & 9.31e3 & 8729.62 & 3184.88 & 2036.51 & 2.8816\\
PPHE3 & 0.6119 & 0.1831 & 7.42e3 & 1.58e4 & 6074.59 & 1787.29 & 1271.25 & 7.5640\\
\bottomrule
\end{tabular}
\end{table}

\section{Limitación pendiente}

Las propiedades físicas de los fluidos del segundo caso de estudio \textbf{no están publicadas
en el artículo}: las toma de su referencia [24],
\begin{quote}
O.P. Arsenyeva, L.L. Tovazhnyanskyy, P.O. Kapustenko, G.L. Khavin, A.P. Yuzbashyan,
P.Y. Arsenyev, \emph{Two types of welded plate heat exchangers for efficient heat recovery in
industry}, Appl. Therm. Eng. 105 (2016) 763--773.
\end{quote}

Mientras tanto el modelo lleva $\rho$ despejada de la Tabla 9 y valores genéricos de
hidrocarburo para $\mu$, $k$ y $c_p$. El script avisa de ello al ejecutarse, mediante la
constante \texttt{PROPIEDADES\_DE\_REF\_24}.

\textbf{La prueba de validación real será la columna $\Delta P_{int}$}: el artículo satura los
100\,kPa admisibles en las seis variantes, de modo que con las propiedades correctas las seis
deberían dar 100. Actualmente van de 49 a 155\,kPa, y esa dispersión es la medida directa del
error que introducen las propiedades provisionales.

\section{Estado del repositorio}

Rama \texttt{integracion-iteracion-velocidad}. Sin confirmar en el momento de escribir este
informe: las modificaciones de \texttt{formulas.py} y \texttt{geometría.py}, y el fichero nuevo
\texttt{modelo\_antiguo.py}.

\end{document}
